\chapter{Закон эволюции}
\label{ch:evolution}

\section{Общая структура}

Дискретность~$\Lambda$ и пара $(\hl,\htick)$ установлены теоремой~\ref{thm:discreteness}
и гл.~\ref{ch:carrier}.
Глобальное состояние эволюционирует по закону
\[
\Psi^t=\{z(x,t)\}_{x\in\Lambda}
\]
\[
  \Psi^{t+1}=g(\Psi^t),\qquad g=g_{\mathrm{nl}}\circ g_{\mathrm{lin}},
\]
где $g_{\mathrm{lin}}$ --- линейное смешивание по соседям~$N(x)$,
$g_{\mathrm{nl}}$ --- умножение на фазу $e^{i\Phi(z)}$.

\begin{definition}[Дискретный лапласиан]
$(\Delta_4 z)(x)=\sum_{y\in N(x)}(z(y)-z(x))$.
\end{definition}

\section{Ограничения шага}

Норма за такт сохраняется,
\[
  \sum_x|z|^2=\mathrm{const}.
\]
Допустимы локальные унитарные преобразования на связях и действительная фаза~$\Phi$.
Нелинейность --- множитель $e^{i\Phi}$.
При $|z|\to 0$ фазовый отклик~$\Phi$ неограниченно возрастает.
При фиксированной норме меняется аргумент.
Знаменатель импульса насыщается добавкой~$+\varepsilon$.
При больших $|z|$ нелинейность затухает.
Фаза шага есть инвариант~$\zeta$.

Линейное выражение $z+\gamma\Delta_4 z$ не унитарно;
допустимо произведение локальных унитарных преобразований.

\section{Фазовый инвариант \texorpdfstring{$\zeta$}{zeta}}

Разность $\Arg(\sum z/z)$ и функция $\atan$ дают на решётке ложный скачок $2\pi$.
Физически значим величина
\[
  \zeta=(\textstyle\sum_N z)\cdot z^*,\qquad \Delta\phi=\wrap(\Arg\zeta),
\]
служащая входом насыщающего фазового импульса.

\section{Спин и группа \texorpdfstring{$SU(2)$}{SU(2)}}

\begin{definition}[Спинорное поле]
$z(x,t)=(z_1,z_2)^\top\in\mathbb{C}^2$, $|z|^2=z^\dagger z$.
\end{definition}

Глобальная $U(1)$-фаза $z\to z e^{i\Phi}$ описывает сохранение заряда,
но недостаточна для фермиона: оборот на $2\pi$ должен менять знак состояния,
на $4\pi$ --- возвращать его.

\begin{definition}[Дискретный вращатель]
\label{def:rotation-gate}
$N_{\mathrm{ring}}=512$, $\omega=e^{i2\pi/N_{\mathrm{ring}}}$, $\Phi\in\mathbb{Z}_{N_{\mathrm{ring}}}$:
\[
  R(\Phi):\ (U,V)\mapsto (U c_\Phi - V s_\Phi,\ U s_\Phi + V c_\Phi)\pmod N,
\]
где $(c_\Phi,s_\Phi)$ --- дискретные cos/sin угла по модулю~$N_{\mathrm{ring}}$.
\end{definition}

\begin{corollary}[Спин $1/2$]
$R(N_{\mathrm{ring}}/2)\mapsto -z$ на первичном вихре $\Rightarrow$
$U(2\pi)=-\mathbb{1}$, $U(4\pi)=+\mathbb{1}$.
\end{corollary}

\begin{proof}
По \cref{def:rotation-gate} угол $\Phi=N_{\mathrm{ring}}/2$ соответствует повороту на~$\pi$
в дискретном $SU(2)$.
На спиноре $(z_1,z_2)$ это даёт $(-z_1,-z_2)=-z$:
\[
  U(2\pi)=-\mathbb{1},\qquad U(4\pi)=+\mathbb{1}.
\]
\end{proof}

Запрет Паули: два параллельных спинора в одной ячейке~$\PlanckCell$
приводят к $z_A+z_B\to 0$; коэффициент~$K_P$ задаёт отталкивание при $d\le\hl$.

\section{$U(1)_{\mathrm{vac}}$ и хиральность}

Глобальная фаза $z\to z e^{i\theta}$; инварианты --- голономия~$\zeta$
и ось Блоха.
\[
n=z^\dagger\sigma z/|z|^2
\]
Проекторы $P_L z=(z_1,0)^\top$, $P_R z=(0,z_2)^\top$
независимы в фазе свободного течения и смешиваются при столкновении
через~$K_P$ и аргумент~$\Arg$.

\section{Замыкание \texorpdfstring{$g$}{g} на \texorpdfstring{$\mathbb{Z}_N[i]$}{Z\_N[i]}}
\label{sec:g-law}

\begin{theorem}[Замыкание закона эволюции]
\label{thm:g-law}
Закон~$g$ на $Z_{N_{\mathrm{ring}}}[i]$:
\begin{enumerate}
\item схема «лягушка'' второго порядка
\[
  Z(x,t+\Delta t)+Z(x,t-\Delta t)=2Z(x,t)+\lfloor\mathcal{N}\rfloor;
\]
\item действительная и мнимая части голономия из суммы соседей без $\atan$;
\item насыщающий фазовый импульс
\[
  \Phi=\Bigl\lfloor\frac{K_P\zeta_{\mathrm{imag}}}{\zeta_{\mathrm{real}}+|Z|^2+K_P}\Bigr\rfloor,\quad
  |\Phi|\ge\Delta\phi_{\mathrm{disc}}\ \text{или}\ \Phi=0;
\]
\item $\lfloor\mathcal{N}\rfloor=R(\Phi)-Z$ на обеих компонентах.
\end{enumerate}
\end{theorem}

\begin{proof}
\emph{Схема «лягушка''.}
Норма $\sum_x|z|^2$ за такт сохраняется.
Обновление первого порядка $z'=z+\gamma\Delta z$ с вещественным~$\gamma$ его не сохраняет:
\[
  |z'|^2-|z|^2=2\gamma\,\Re(z^*\Delta z)+\gamma^2|\Delta z|^2,
\]
и правая часть не равна нулю тождественно.
Допустима форма второго порядка по времени, в которой шаг вперёд и шаг назад входят симметрично:
\[
  Z(x,t+\Delta t)+Z(x,t-\Delta t)=2Z(x,t)+\lfloor\mathcal{N}(x,t)\rfloor.
\]
Неизвестный предыдущий регистр выражается явно,
\[
  Z(x,t-\Delta t)=2Z(x,t)+\lfloor\mathcal{N}\rfloor-Z(x,t+\Delta t),
\]
и это есть $g^{-1}$.
Округление $\lfloor\cdot\rfloor$ --- работа в кольце $\mathbb{Z}_{N_{\mathrm{ring}}}[i]$.

\smallskip
\emph{Голономия.}
Разность $\Arg(z(y)/z(x))$ на решётке скачет на $2\pi$.
Фаза шага есть
\[
  \zeta=\Bigl(\sum_{y\in N(x)} z(y)\Bigr)\cdot z(x)^*.
\]
В удар входят $\zeta_{\mathrm{real}}$ и $\zeta_{\mathrm{imag}}$.

\smallskip
\emph{Насыщающий импульс.}
При $|z|\to 0$ знаменатель обращается в нуль, и фазовый отклик $\Phi$ неограничен.
При больших $|z|$ рост ограничен добавкой $|Z|^2+K_P$ в знаменателе:
\[
  \Phi=\Bigl\lfloor\frac{K_P\,\zeta_{\mathrm{imag}}}{\zeta_{\mathrm{real}}+|Z|^2+K_P}\Bigr\rfloor,
\]
\[
  \Phi=0\quad\text{или}\quad |\Phi|\ge\Delta\phi_{\mathrm{disc}}.
\]

\smallskip
\emph{Вращатель.}
Нелинейность шага есть унитарная фаза на~$\mathbb{C}^2$:
\[
  \lfloor\mathcal{N}\rfloor=R(\Phi)-Z
\]
на обеих компонентах (\cref{def:rotation-gate}).
\end{proof}

\subsection{Информационный бюджет ячейки}
\label{sec:bit-budget}

Нативная запись бюджета и координат $(dh,dt)$, $(k_\phi,\varphi_{\mathrm{f}},\mu)$
--- \cref{ch:native-coords,def:dh-dt,def:internal-coords,rem:phi-omega-Phi}.

\[
  B_{\hv}=\frac{2\pi}{\ln 2}\approx 9{,}065\ \text{бит},\quad
  N_{\mathrm{ring}}=2^{\lfloor B_{\hv}\rfloor}=512,\quad
  N_\phi=\lceil 4\pi\rceil=13,\quad
  \Delta\phi_{\mathrm{disc}}=41.
\]

\subsection{Лестница согласованности}
\label{sec:congruence}

\begin{proposition}[CL-1 -- CL-8]
По модулю $N_{\mathrm{ring}}=512$ шаг обратим, половина кольца меняет знак спинора, энергия и импульс кратны кванту, а $41$ и $13$ порождают кольцо.
\end{proposition}

\begin{proof}
\emph{CL-1.}
Из схемы второго порядка
\[
  Z^++Z^-=2Z+\lfloor\mathcal{N}\rfloor
\]
предыдущий регистр выражается явно:
\[
  Z^-=2Z+\lfloor\mathcal{N}\rfloor-Z^+.
\]

\emph{CL-2.}
Шаг кольца есть $2\pi/N_{\mathrm{ring}}$, поэтому половина кольца --- угол $\pi$:
\[
  \frac{2\pi}{N_{\mathrm{ring}}}\cdot\frac{N_{\mathrm{ring}}}{2}=\pi,
  \qquad
  R(N_{\mathrm{ring}}/2)\,z=-z.
\]

\emph{CL-3.}
\[
  \Phi=0\quad\text{или}\quad |\Phi|\ge\Delta\phi_{\mathrm{disc}}.
\]

\emph{CL-4.}
\[
  n_E=\Bigl\lfloor\frac{|\Phi|}{\Delta\phi_{\mathrm{disc}}}\Bigr\rfloor,
  \qquad
  E=n_E E_0.
\]

\emph{CL-5, CL-6.}
\[
  N_{\mathrm{ring}}=2^9=512,
  \qquad
  \Delta\phi_{\mathrm{disc}}=41,
  \qquad
  N_\phi=13,
\]
\[
  \gcd(41,512)=1,
  \qquad
  \gcd(13,512)=1.
\]
Кратности $41$ пробегают всё кольцо $\mathbb{Z}/512\mathbb{Z}$.
Умножение на $13$ в этом кольце обратимо, тринадцать секторов различны.

\emph{CL-7, CL-8.}
\[
  \Delta E=n_E E_0,
  \qquad
  \Delta p=n_p p_0,
  \qquad
  n_E,n_p\in\mathbb{Z}.
\]
На замкнутом шаге
\[
  \sum\Delta E=E_0\sum n_E,
  \qquad
  \sum\Delta p=p_0\sum n_p.
\]
\end{proof}

Топологический заряд $n\in\mathbb{Z}$. Регистр фазы --- вычет по модулю~$N_{\mathrm{ring}}$.
